\section*{PHẦN MỞ ĐẦU}
\addcontentsline{toc}{section}{PHẦN MỞ ĐẦU}

\subsection*{1. Lý do chọn đề tài}
\addcontentsline{toc}{subsection}{1. Lý do chọn đề tài}

Trong toàn bộ hệ thống lý luận của chủ nghĩa duy vật lịch sử, mối quan hệ biện chứng giữa tồn tại xã hội và ý thức xã hội, giữa cơ sở hạ tầng và kiến trúc thượng tầng luôn giữ vị trí nguyên lý nền tảng để nhận thức tiến trình vận động của lịch sử nhân loại. Triết học Mác -- Lênin khẳng định rằng phương thức sản xuất vật chất quyết định diện mạo tinh thần của xã hội; đồng thời ý thức xã hội có tính độc lập tương đối, có khả năng tác động trở lại mạnh mẽ đối với sự phát triển của đời sống hiện thực. Trong cấu trúc xã hội trung đại Việt Nam, tầng lớp nho sĩ giữ vai trò là chủ thể trực tiếp mang tải, định hình và thực thi hệ tư tưởng chính thống, chi phối sâu sắc đời sống chính trị, văn hóa và đạo đức của quốc gia.

Thế kỷ XVIII ở Việt Nam là một giai đoạn lịch sử đặc biệt, chứng kiến những cuộc chuyển mình to lớn và các cơn biến động dữ dội nhất trong kỷ nguyên phong kiến. Đây là thời kỳ chế độ quân chủ chuyên chế lâm vào cuộc khủng hoảng trầm trọng ở cả Đàng Ngoài lẫn Đàng Trong, làm suy sụp uy tín và vị thế độc tôn của hệ tư tưởng Nho giáo vốn được xác lập từ thời Lê sơ. Cùng với sự mục ruỗng của thể chế phong kiến cũ, cao trào khởi nghĩa nông dân bùng nổ mạnh mẽ khắp mọi miền, mà đỉnh cao là phong trào Tây Sơn do người anh hùng dân tộc Quang Trung -- Nguyễn Huệ lãnh đạo, đã đập tan các tập đoàn phong kiến cát cứ, quét sạch quân xâm lược Xiêm và Mãn Thanh, đặt nền móng vững chắc cho sự nghiệp thống nhất giang sơn.

Đứng trước sự biến đổi sâu sắc của tồn tại xã hội, tầng lớp nho sĩ Việt Nam thế kỷ XVIII không còn duy trì được sự thuần nhất về mặt tư tưởng. Họ trải qua một quá trình phân hóa sâu rộng về lập trường chính trị và đạo đức xử thế: từ xu hướng bảo thủ cố chấp níu giữ ngai vàng phong kiến đã suy tàn, đến thái độ bất mãn lánh đời mở trường tư thục, cho đến đỉnh cao là sự thức thời dấn thân đồng hành cùng phong trào cách mạng nông dân Tây Sơn để phụng sự Tổ quốc. Không dừng lại ở đó, thế kỷ XVIII còn đánh dấu bước chuyển mang tính đột phá trong tư duy học thuật: sự suy thoái của lối học kinh viện Tống Nho giáo điều đã mở đường cho trào lưu thực học, tư duy bách khoa khảo chứng, xu hướng hòa đồng Tam giáo và sự bùng nổ của chủ nghĩa nhân đạo sâu sắc trong các trước tác văn học chữ Nôm.

Việc nghiên cứu đề tài ``Tìm hiểu về tầng lớp nho sĩ ở Việt Nam thế kỷ XVIII'' mang ý nghĩa cấp thiết trên cả hai phương diện:

Về phương diện lý luận, đề tài góp phần làm sáng tỏ quy luật duy vật lịch sử về sự tác động của cơ sở kinh tế, xã hội đối với sự chuyển hóa của ý thức hệ; phân tích bản chất hai mặt, vị trí xã hội và quy luật phân hóa của tầng lớp trí thức phong kiến trong những bước ngoặt trọng đại của dân tộc.

Về phương diện thực tiễn, việc tổng kết sự chuyển biến của nho sĩ thế kỷ XVIII đem lại những bài học kinh nghiệm sâu sắc cho giai đoạn hiện nay: bài học về nguyên tắc thống nhất giữa lý luận với thực tiễn, kiên quyết bài trừ thói chuộng danh ảo và bệnh bằng cấp hình thức; bài học về việc rèn luyện bản lĩnh chính trị của người trí thức trước những thử thách của thời cuộc; và bài học về chính sách chiêu hiền đãi sĩ, trọng dụng nhân tài nhằm phát huy cao độ nguồn nguyên khí quốc gia trong công cuộc xây dựng và bảo vệ Tổ quốc Việt Nam xã hội chủ nghĩa.

\subsection*{2. Tổng quan tình hình nghiên cứu}
\addcontentsline{toc}{subsection}{2. Tổng quan tình hình nghiên cứu}

Vấn đề lịch sử, tư tưởng và diện mạo tầng lớp nho sĩ Việt Nam thế kỷ XVIII là đối tượng thu hút sự quan tâm nghiên cứu của nhiều học giả chuyên ngành triết học, sử học và văn hóa học. Các công trình tiêu biểu đã công bố có thể khái quát thành ba hướng tiếp cận chủ yếu:

Thứ nhất, nhóm các công trình nghiên cứu tổng quan về lịch sử tư tưởng và triết học Việt Nam. Điển hình là công trình \textit{Lịch sử tư tưởng Việt Nam} (Tập 1) do Giáo sư, Tiến sĩ Nguyễn Tài Thư chủ biên (Viện Triết học, Nhà xuất bản Khoa học Xã hội, 1993) đã phân tích sâu sắc quá trình suy thoái của Nho giáo chính thống từ thế kỷ XVI đến thế kỷ XVIII, làm nổi bật sự xuất hiện của các mầm mống tư tưởng tiến bộ và trào lưu thực học. Cùng với đó, tác phẩm \textit{Sự phát triển của tư tưởng ở Việt Nam từ thế kỷ XIX đến Cách mạng Tháng Tám} (Tập 1) của Giáo sư Trần Văn Giàu (Nhà xuất bản Khoa học Xã hội, 1973) đã làm rõ bối cảnh khủng hoảng ý thức hệ phong kiến cuối thế kỷ XVIII như một tiền đề tư tưởng trực tiếp dẫn tới những biến chuyển xã hội ở các thế kỷ tiếp theo.

Thứ hai, nhóm các công trình nghiên cứu lịch sử chế độ phong kiến và văn hóa giáo dục trung đại. Bộ sách \textit{Lịch sử chế độ phong kiến Việt Nam} (Tập III) của tập thể tác giả Phan Huy Lê, Chu Thiên, Vương Hoàng Tuyên, Đinh Xuân Lâm (Nhà xuất bản Giáo dục, 1965) đã tái hiện chân thực bức tranh khủng hoảng giai cấp, sự suy đồi của quan trường và phong trào đấu tranh của nông dân thế kỷ XVIII. Cùng với hướng nghiên cứu này, tác phẩm \textit{Việt Nam văn hóa sử cương} của học giả Đào Duy Anh và bộ sách \textit{Đại cương lịch sử văn hóa Việt Nam} (Tập 4) của nhà nghiên cứu Nguyễn Khắc Thuần đã cung cấp những khảo sát cụ thể về chế độ thi cử, việc học tập và hiện tượng tha hóa khoa cử ``Sinh đồ ba quan'' năm 1750 thời Lê -- Trịnh.

Thứ ba, nhóm các bài báo khoa học và chuyên khảo chuyên sâu về danh nhân, giáo dục và tư tưởng thế kỷ XVIII. Bài báo của tác giả Nguyễn Thị Hồng Phương đăng trên \textit{Tạp chí Nghiên cứu Triết học} (2026) cùng các công trình của Viện Triết học đã đi sâu phân tích bản thể luận về phạm trù ``Khí'', nhận thức luận ``Lý ở trong Khí'' và nhân sinh quan nhập thế của nhà bác học Lê Quý Đôn; bên cạnh việc khảo cứu tư tưởng hòa đồng Tam giáo của Ngô Thì Nhậm qua tác phẩm \textit{Trúc Lâm tông chỉ nguyên thanh}. Ngoài ra, công trình của Nguyễn Kim Sơn về những nhân tố khởi phát khuynh hướng khảo chứng học, bài viết của Đặng Việt Bích về hình ảnh vua Quang Trung trong tâm trí kẻ sĩ Bắc Hà, nghiên cứu của Quách Hiền về phong trào tư học, khảo cứu của Cao Văn Thức về học vị Sinh đồ thời Lê mạt và chuyên đề tư liệu của Sở Giáo dục và Đào tạo Thành phố Hồ Chí Minh cùng Thư viện tỉnh Đồng Nai về danh sư Võ Trường Toản đã cung cấp nguồn dữ liệu thực chứng quý báu.

Kế thừa có chọn lọc thành tựu của các học giả đi trước, tiểu luận này đi sâu khảo sát diện mạo tầng lớp nho sĩ thế kỷ XVIII một cách có hệ thống dưới lăng kính phương pháp luận của chủ nghĩa duy vật biện chứng và chủ nghĩa duy vật lịch sử, kết hợp chặt chẽ giữa phân tích thái độ chính trị với việc đánh giá bước nhảy vọt trong tư tưởng học thuật, từ đó đúc kết những bài học có giá trị chỉ đạo thực tiễn sâu sắc.

\subsection*{3. Mục tiêu và nhiệm vụ nghiên cứu}
\addcontentsline{toc}{subsection}{3. Mục tiêu và nhiệm vụ nghiên cứu}

Mục tiêu tổng quát của đề tài là vận dụng thế giới quan và phương pháp luận mácxít nhằm làm sáng tỏ diện mạo, sự phân hóa về lập trường chính trị -- xã hội và những bước chuyển mình sâu sắc trong tư tưởng triết học, khuynh hướng học thuật của tầng lớp nho sĩ Việt Nam thế kỷ XVIII; qua đó đánh giá khách quan các đóng góp tiến bộ cùng hạn chế lịch sử, rút ra bài học kinh nghiệm thiết thực cho công tác xây dựng đội ngũ trí thức nước nhà trong giai đoạn phát triển hiện nay.

Để hoàn thành mục tiêu tổng quát nêu trên, đề tài xác định và giải quyết ba nhiệm vụ cụ thể tương ứng với ba chương của công trình:

\textit{Nhiệm vụ thứ nhất,} làm rõ cơ sở lý luận triết học Mác -- Lênin về mối quan hệ biện chứng giữa tồn tại xã hội với ý thức xã hội, đặc điểm của tầng lớp trí thức trong xã hội có giai cấp; đồng thời phân tích các điều kiện lịch sử, kinh tế, chính trị và giáo dục khoa cử của Việt Nam thế kỷ XVIII đã tác động trực tiếp đến sự biến đổi của tầng lớp nho sĩ.

\textit{Nhiệm vụ thứ hai,} phân tích thực trạng phân hóa về thái độ chính trị -- xã hội của tầng lớp nho sĩ trước các biến động thời cuộc; đồng thời làm rõ những chuyển biến lớn trong thế giới quan triết học, trào lưu thực học bách khoa, sự dung hợp Tam giáo và ý thức nhân đạo cao cả của nho sĩ thế kỷ XVIII.

\textit{Nhiệm vụ thứ ba,} đánh giá toàn diện những giá trị đóng góp tiến bộ và hạn chế mang tính lịch sử của tầng lớp nho sĩ thế kỷ XVIII dưới góc độ duy vật lịch sử; từ đó đúc kết những bài học kinh nghiệm có ý nghĩa khoa học và thực tiễn đối với sự nghiệp xây dựng, phát huy vai trò của đội ngũ trí thức Việt Nam trong kỷ nguyên vươn mình của dân tộc.

\subsection*{4. Đối tượng và phạm vi nghiên cứu}
\addcontentsline{toc}{subsection}{4. Đối tượng và phạm vi nghiên cứu}

Đối tượng nghiên cứu của đề tài là tầng lớp nho sĩ ở Việt Nam thế kỷ XVIII, trọng tâm là sự phân hóa về lập trường chính trị -- xã hội và sự vận động, chuyển biến trong thế giới quan triết học, nhân sinh quan và phương pháp học thuật của họ dưới sự quy định của hoàn cảnh lịch sử khách quan.

Phạm vi nghiên cứu được xác định cụ thể trên hai bình diện:

Về không gian, đề tài khảo sát trên toàn bộ lãnh thổ đất nước Đại Việt thế kỷ XVIII, bao quát cả hai xứ Đàng Ngoài và Đàng Trong.

Về thời gian, đề tài tập trung nghiên cứu trong phạm vi thế kỷ XVIII, từ giai đoạn khủng hoảng trầm trọng của chế độ phong kiến Lê -- Trịnh và chúa Nguyễn đến thời kỳ bão táp của phong trào khởi nghĩa Tây Sơn cuối thế kỷ XVIII; đồng thời có sự liên hệ với tiền đề Nho giáo thời Lê sơ thế kỷ XV và mở rộng liên hệ với thực tiễn xây dựng đội ngũ trí thức hiện nay.

\subsection*{5. Phương pháp nghiên cứu}
\addcontentsline{toc}{subsection}{5. Phương pháp nghiên cứu}

Đề tài này được thực hiện trên cơ sở phương pháp luận của chủ nghĩa duy vật biện chứng và chủ nghĩa duy vật lịch sử. Ngoài ra, đề tài còn sử dụng các phương pháp khác như:

Phương pháp thống nhất giữa lịch sử và logic nhằm tái hiện tiến trình vận động cụ thể của lịch sử dân tộc thế kỷ XVIII, qua đó khái quát bản chất và quy luật biến đổi tư tưởng của tầng lớp nho sĩ.

Phương pháp phân tích và tổng hợp được vận dụng để mổ xẻ từng khuynh hướng thái độ, từng trào lưu học thuật của các nhóm nho sĩ, sau đó tổng hợp thành diện mạo tổng thể về tầng lớp trí thức thế kỷ XVIII.

Phương pháp so sánh, đối chiếu được sử dụng nhằm làm rõ sự khác biệt giữa các xu hướng hành xử bảo thủ và tiến bộ, giữa nho sĩ Bắc Hà và nho sĩ Nam Bộ, giữa lối học kinh viện Tống Nho và trào lưu thực học bách khoa.

Phương pháp tiếp cận hệ thống đặt tầng lớp nho sĩ trong mối quan hệ biện chứng với cơ sở kinh tế, thể chế chính trị và phong trào đấu tranh cách mạng của quần chúng nhân dân đương thời.

\subsection*{6. Kết cấu của đề tài}
\addcontentsline{toc}{subsection}{6. Kết cấu của đề tài}

Ngoài phần mở đầu, kết luận, tài liệu tham khảo, nội dung của đề tài này bao gồm 03 Chương, 06 Tiết. Ba chương được triển khai theo mạch logic chặt chẽ từ việc xác lập cơ sở lý luận duy vật lịch sử và phân tích bối cảnh biến động kinh tế -- xã hội thế kỷ XVIII, tiến tới mổ xẻ thực trạng phân hóa lập trường chính trị cùng bước chuyển biến tư tưởng học thuật của tầng lớp nho sĩ, qua đó đánh giá khách quan các giá trị, hạn chế lịch sử và đúc rút bài học kinh nghiệm đối với việc xây dựng đội ngũ trí thức Việt Nam hiện nay.

